%! Piecewise & Step Functions — Unit 1, Lesson 1.4 — Homework
% GRADUAL RELEASE (stats/SAAR shape, adopted for 1.4 on 2026-09-07): this page IS the lesson's
% individual practice. It is scored, it is the last component in the packet, and it is STARTED
% IN CLASS in the last ten minutes of the period, alone and silent, then finished at home. The
% teacher may substitute a DeltaMath set on a given day; that replaces this page and strikes its
% score cell on the cover. Due the first class after two study halls — never "next class".
% TWO PAGES IS A CEILING at 12pt: two boxes — items 1-3 on page 1, items 4-6 and the spiralbox
% on page 2 — so no item breaks across a page. NO EXTENSION BOX (retired with this shape).
% Six items span the whole of A2.F.2a/b/c/f: the core procedure off a rule (1), the
% meets-or-jumps contrast pair (2), the same skill off a graph — read it AND write its rule (3),
% the special case and its boundary (4), a model with an interpret-the-answer follow-up (5),
% and the SOL-style multiple-choice formative check (6).
% Every work block and every \writelines{n} is matched byte-for-byte / n-for-n in the key.
\documentclass[12pt]{article}
\usepackage{algebra2-article}
\usepackage{algebra2-boxes}
% Coordinate grid: {xmin}{xmax}{ymin}{ymax}
\newcommand{\lgrid}[4]{%
  \draw[step=1, linegray!40, very thin] (#1,#3) grid (#2,#4);
  \draw[->, semithick, charcoal] (#1,0)--(#2,0) node[below right, font=\scriptsize]{$x$};
  \draw[->, semithick, charcoal] (0,#3)--(0,#4) node[left, font=\scriptsize]{$y$};}
% Endpoint dots: {color}{x}{y}. Closed = filled; open = white-filled ring.
\newcommand{\fdot}[3]{\fill[#1] (#2,#3) circle (3.2pt);}
\newcommand{\hdot}[3]{\draw[very thick, #1, fill=white] (#2,#3) circle (3.2pt);}

\begin{document}

\pageheader{Unit 1, Lesson 1.4}{Homework}
% NAMESTRIP: no \namedateperiod here — the name row lives on the cover only.

\vspace{0.10in}

\boxguard[8]
\begin{remindbox}
\small
\textbf{This is your graded homework.} Your packet with completed homework is due the first
class after two study halls. I will announce the due date in class and on TurtleNet. Remember:
first find the piece whose condition the input satisfies --- at a \textbf{boundary}, the piece
whose inequality \emph{includes} it --- then apply that rule. A closed dot includes its endpoint;
an open dot excludes it.
\end{remindbox}

\vspace{0.10in}

\boxguard[14]
\begin{notesbox}{Practice}
\small
\begin{enumerate}[leftmargin=1.9em, itemsep=8pt]
  \item \textbf{Evaluate.} $f(x)=\begin{cases} x+4, & x\le -1 \\ 2x, & x>-1 \end{cases}$
    \\[3pt]
    $f(-3)=$ \blank{1.1cm} \quad $f(-1)=$ \blank{1.1cm} \quad $f(0)=$ \blank{1.1cm} \quad
    $f(5)=$ \blank{1.1cm} \\[3pt]
    Which piece owns $x=-1$, and why? \blank{6.0cm}

  \item \textbf{Meets or jumps?} Two functions, one number apart:
    \[ g(x)=\begin{cases} x+2, & x<1 \\ 5-x, & x\ge1 \end{cases} \qquad\qquad
       k(x)=\begin{cases} x+2, & x<1 \\ 4-x, & x\ge1 \end{cases} \]
    \vspace{-8pt}
    \begin{center}
    \renewcommand{\arraystretch}{1.4}
    \begin{tabularx}{\linewidth}{@{}X|X@{}}
    \textbf{(a) $g$} & \textbf{(b) $k$} \\
    \hline
    At $x=1$: first rule \blank{0.8cm}, second rule \blank{0.8cm} & At $x=1$: first rule \blank{0.8cm}, second rule \blank{0.8cm} \\
    So the graph \blank{1.2cm} (meets / jumps) & So the graph \blank{1.2cm} (meets / jumps) \\
    \end{tabularx}
    \end{center}
    (c) In one sentence: what decides whether two pieces meet at a boundary? \writelines{2}

  \item \textbf{Read the graph, then write its rule.} The graph shows \emph{all} of $f$. \\[2pt]
    \begin{minipage}[c]{0.60\linewidth}
    $f(1)=$ \blank{0.8cm}; \; $f(2)=$ \blank{0.8cm} (read the \blank{1.0cm} dot) \\[3pt]
    Domain \blank{1.6cm}; \; increasing on \blank{1.6cm}; \; constant on \blank{1.6cm} \\[3pt]
    Continuous at $x=2$? \blank{0.9cm} \; Now write the rule: \\[3pt]
    $f(x)=\begin{cases} \text{\blank{1.2cm}}, & -2\le x<2 \\[3pt] \text{\blank{1.2cm}}, & 2\le x\le5 \end{cases}$
    \end{minipage}\hfill
    \begin{minipage}[c]{0.36\linewidth}\centering
    \begin{tikzpicture}[scale=0.38]
      \lgrid{-2.5}{5.5}{-3}{3}
      \draw[very thick, forest] (-2,-2)--(2,2);
      \fdot{forest}{-2}{-2} \hdot{forest}{2}{2}
      \draw[very thick, forest] (2,-2)--(5,-2);
      \fdot{forest}{2}{-2} \fdot{forest}{5}{-2}
    \end{tikzpicture}
    \end{minipage}
\end{enumerate}
\end{notesbox}

\newpage
\begin{notesbox}{Practice, continued}
\small
\begin{enumerate}[leftmargin=1.9em, itemsep=8pt, start=4]
  \item \textbf{The staircase.} Round \emph{down} each time.
    \\[3pt]
    $\lfloor 4.2\rfloor=$ \blank{0.9cm} \quad $\lfloor -3.1\rfloor=$ \blank{0.9cm} \quad
    $\lfloor 6\rfloor=$ \blank{0.9cm} \quad $\lfloor -0.5\rfloor=$ \blank{0.9cm} \quad
    For every $x$ in $[3,4)$, $\lfloor x\rfloor=$ \blank{0.9cm} \\[3pt]
    For which inputs is $\lfloor x\rfloor=x$ exactly? \blank{2.6cm} \; Is $\lfloor x\rfloor$
    \emph{ever} bigger than $x$? \blank{0.9cm}

  \item \textbf{Model it.} A phone plan costs \$30 for up to $5$ GB of data, then \$10 for each
        GB over $5$. Let $C(g)$ be the monthly cost for $g$ GB.
    \begin{enumerate}[label=(\alph*), leftmargin=1.9em, itemsep=4pt]
      \item Complete the rule:
            $C(g)=\begin{cases} \text{\blank{1.4cm}}, & 0\le g\le5 \\[3pt] 30+\text{\blank{2.4cm}}, & g>5 \end{cases}$
      \item Find $C(8)$, and say what it means.
        \begin{work}
          C(8) &= 30 + 10(8-5) \\
               &= 60
        \end{work}
        $C(8)$ means: \blank{7.6cm}
      \item Now the company adds a \$10 \emph{surcharge} the moment you pass $5$ GB. Rewrite the
            second piece: $C(g)=$ \blank{3.0cm} for $g>5$. Do the pieces still meet at $g=5$?
            \blank{0.9cm} \; The graph now \blank{1.2cm} from \$\blank{0.8cm} to \$\blank{0.8cm}.
    \end{enumerate}

  \item \textbf{Multiple choice (SOL-style).} For
        $f(x)=\begin{cases} 2x+1, & x<2 \\ x-4, & x\ge2 \end{cases}$, what is $f(2)$?
    \begin{enumerate}[label=\Alph*., leftmargin=2.2em, itemsep=1pt]
      \item $5$
      \item $-2$
      \item $3$
      \item $f(2)$ is undefined
    \end{enumerate}
    Say in one sentence how you ruled out one of the others. \writelines{2}
\end{enumerate}
\end{notesbox}

\vspace{0.10in}

\boxguard[8]
\begin{spiralbox}
\small
\textbf{Coming next --- Lesson 1.5: Linear Regression.} We leave exact rules behind. Given a
\emph{scatter} of real data points that lie on no single line, you will read one \textbf{line of
best fit}, measure how well it fits with the \textbf{correlation coefficient}, and use it to
predict --- the first time in Unit~1 the data is messy.
\end{spiralbox}

\end{document}
